\input{\commonfolder/header}

%% Chinese versions of the user-visible common macros.
\renewcommand{\seedsubmission}{你需要提交一份详细的实验报告，并附上截图，
用来说明你完成了什么、观察到了什么。对于有趣或出乎意料的观察结果，
你还需要给出解释。请同时列出重要的代码片段，并对这些代码进行解释。
只提交代码而没有解释将不会获得分数。}

\renewcommand{\seedbook}{\textit{Computer \& Internet Security: A Hands-on Approach}，
第 3 版，Wenliang Du 著。详情见 \url{https://www.handsonsecurity.net}。\xspace}

\renewcommand{\seedisbook}{\textit{Internet Security: A Hands-on Approach}，
第 3 版，Wenliang Du 著。详情见 \url{https://www.handsonsecurity.net}。\xspace}

\renewcommand{\seedcsbook}{\textit{Computer Security: A Hands-on Approach}，
第 3 版，Wenliang Du 著。详情见 \url{https://www.handsonsecurity.net}。\xspace}

\renewcommand{\seedcibook}{\textit{Computer \& Internet Security: A Hands-on Approach}，
第 3 版，Wenliang Du 著。详情见 \url{https://www.handsonsecurity.net}。\xspace}

\renewcommand{\seedisvideo}{\textit{Internet Security: A Hands-on Approach}，
Wenliang Du 主讲。详情见 \url{https://www.handsonsecurity.net/video.html}。\xspace}

\renewcommand{\seedcsvideo}{\textit{Computer Security: A Hands-on Approach}，
Wenliang Du 主讲。详情见 \url{https://www.handsonsecurity.net/video.html}。\xspace}

\renewcommand{\seedenvironment}{本实验已经在我们预先构建的 Ubuntu 16.04
虚拟机上测试过。该虚拟机可以从 SEED 网站下载。\xspace}

\renewcommand{\seedenvironmentA}{本实验已经在我们预先构建的 Ubuntu 16.04
虚拟机上测试过。该虚拟机可以从 SEED 网站下载。\xspace}

\renewcommand{\seedenvironmentB}{本实验已经在我们预先构建的 Ubuntu 20.04
虚拟机上测试过。该虚拟机可以从 SEED 网站下载。\xspace}

\renewcommand{\seedenvironmentC}{本实验已经在 SEED Ubuntu 20.04 虚拟机上
测试过。你可以从 SEED 网站下载预先构建的镜像，并在自己的计算机上运行。
不过，大多数 SEED 实验也可以在云端完成；你可以按照我们的说明在云端创建
SEED 虚拟机。\xspace}

\renewcommand{\seedenvironmentAB}{本实验已经在我们预先构建的 Ubuntu 16.04
和 20.04 虚拟机上测试过。虚拟机可以从 SEED 网站下载。\xspace}

\renewcommand{\nodependency}{由于我们使用容器搭建实验环境，本实验对 SEED
虚拟机的依赖并不强。你也可以在其他虚拟机、物理机或云端虚拟机上完成本实验。\xspace}

\renewcommand{\adddns}{你需要把所需的 IP 地址映射加入
\texttt{/etc/hosts} 文件。\xspace}
